\studytypesubsection{LLMs for New Software Engineering Tools}
\label{sec:llms-for-new-software-engineering-tools}

In this role, LLMs serve as components of new tools that assist software engineers (e.g., with code comprehension or test generation) or as autonomous agents that perform multi-step tasks on their behalf.

\studytypeparagraph{Description}

New LLM-based tools support software engineers in their daily tasks, such as code comprehension~\cite{DBLP:conf/chi/YanHWH24} and test case generation~\cite{DBLP:journals/tse/SchaferNET24}.
% GROUNDING DBLP:conf/chi/YanHWH24: "Ivie improved understanding of generated code, and was received by programmers as a highly useful, low distraction complement to the programming assistant."
% GROUNDING DBLP:journals/tse/SchaferNET24: "We implement our approach in TestPilot, an adaptive LLM-based test generation tool for JavaScript that automatically generates unit tests for the methods in a given project's API."
One way of integrating LLM-based tools into software engineers' workflows is using GenAI agents.
Unlike traditional LLM-based tools, these agents are capable of acting autonomously and proactively, are often tailored to meet specific user needs (e.g., via context files or domain-specific tools), and can interact with external environments (e.g., file systems, shells, or web APIs)~\cite{wiesinger2025agents,yang2025code}.
% GROUNDING wiesinger2025agents: "Agents are autonomous and can act independently of human intervention... Agents can also be proactive in their approach to reaching their goals."
% GROUNDING yang2025code: "the emergence of autonomous coding agents capable of multi-step problem decomposition"
From an architectural perspective, GenAI agents can be implemented in various ways~\cite{wiesinger2025agents}, but at their core they run a control loop around the LLM (observe $\rightarrow$ inspect $\rightarrow$ choose $\rightarrow$ act)~\cite{raschka2026components}.
% GROUNDING wiesinger2025agents: "there are many such architectures that can be achieved by the mixing and matching of these components"
% GROUNDING raschka2026components: "An agent is a layer on top, which can be understood as a control loop around the model."
% GROUNDING raschka2026components: "the agent is the system that repeatedly calls the model inside an environment"
Each chosen action (e.g., a tool call) produces a result that the agent feeds back into the next iteration, until the task is complete.
\emph{CoALA}~\cite{DBLP:journals/tmlr/SumersYN024} offers a conceptual framework for organizing such agents.
% GROUNDING DBLP:journals/tmlr/SumersYN024: "modular memory components, a structured action space to interact with internal memory and external environments, and a generalized decision-making process to choose actions"
For coding agents specifically, \citeauthor{raschka2026components} identifies building blocks such as the repository context gathered before each call, the constructed prompt and its tool definitions, structured session memory, and delegation to bounded subagents~\cite{raschka2026components}.
% GROUNDING raschka2026components: "In this article, I lay out six of the main building blocks of a coding agent."
Because these architectures vary, researchers can test and compare them to study how design choices affect downstream performance.

\studytypeparagraph{Examples}

\citeauthor{DBLP:conf/chi/YanHWH24} proposed \emph{IVIE}, a tool integrated into the VS Code graphical interface that generates and explains code using LLMs~\cite{DBLP:conf/chi/YanHWH24}.
% GROUNDING DBLP:conf/chi/YanHWH24: "We present an implementation of Ivie that forks VS Code, applying a modern LLM for timely segmentation and explanation of generated code."
The authors focused more on the presentation, providing a user-friendly interface to interact with the LLM. 
\citeauthor{DBLP:journals/tse/SchaferNET24}~\cite{DBLP:journals/tse/SchaferNET24} presented a large-scale empirical evaluation on the effectiveness of LLMs for automated unit test generation.
% GROUNDING DBLP:journals/tse/SchaferNET24: "This paper presents a large-scale empirical evaluation on the effectiveness of LLMs for automated unit test generation without requiring additional training or manual effort."
They presented \emph{TestPilot}, a tool that implements an approach in which the LLM is provided with prompts that include the signature and implementation of a function under test, along with usage examples extracted from the documentation.
\citeauthor{DBLP:conf/icsm/RichardsW24} introduced a preliminary GenAI agent designed to assist developers in understanding source code by incorporating a reasoning component grounded in the theory of mind~\cite{DBLP:conf/icsm/RichardsW24}.
% GROUNDING DBLP:conf/icsm/RichardsW24: "we designed an LLM-based conversational assistant that provides a personalized interaction based on inferred user mental state (e.g., background knowledge and experience)"
\citeauthor{DBLP:conf/icse-seip/TakerngsaksiriPTTZJLCCW25}~\cite{DBLP:conf/icse-seip/TakerngsaksiriPTTZJLCCW25} present \emph{HULA}, a multi-agent system deployed in Atlassian JIRA that lets engineers refine LLM-generated coding plans and source code, and report acceptance and modification rates from real users.
% GROUNDING DBLP:conf/icse-seip/TakerngsaksiriPTTZJLCCW25: "We design the user interface, implement, and deploy HULA into Atlassian JIRA for internal uses"
